%% kapa: exam=1 round=35 year=2024 period=1 form=A subject=economics subject_ko=경제학원론 seq=33 number=73
%% source: Q-Net 감정평가사 기출문제 (cst00309) artlSeq=5213229 file=raw/1008/5213229/1교시 문제지 원본.pdf
%% conversion: kordoc 4.24.51 -> tools/gen1.py (auto); review=proofread-sonnet; flags=none
\begin{question}{73}
1인당 생산함수가 $y=0.5k^{0.2}$, 자본의 감가상각률이 0.1, 저축률이 0.2인 솔로우(Solow) 경제성장모형에 관한 설명으로 옳은 것을 모두 고른 것은? (단, $y$는 1인당 생산량, $k$는 1인당 자본량이고, 인구증가와 기술진보는 없다.)
\begin{framedtext}
ㄱ. 현재 1인당 자본량이 2일 때, 1인당 투자는 증가한다. \par
ㄴ. 현재 1인당 자본량이 2일 때, 1인당 자본의 감가상각은 1인당 저축보다 작다. \par
ㄷ. 균제상태(steady state)에서 벗어나 있는 경우, 현재 1인당 자본량에 관계없이, 1인당 생산량의 변화율은 0으로 수렴한다. \par
ㄹ. 균제상태의 1인당 자본량은 황금률(Golden Rule) 수준과 같다. \par
\end{framedtext}
\begin{choices}
  \item ㄱ, ㄴ
  \item ㄱ, ㄷ
  \item ㄴ, ㄷ
  \item ㄴ, ㄹ
  \item ㄷ, ㄹ
\end{choices}
\end{question}
\answer{5}
